\chapter{約束の果実}

\section{物語プロット}
\subsection{シーン１：回想（2年前の庭）}
画面全体が暖かな陽光に包まれている。
まだ元気な母エレナと、8歳の幼いノアが、泥だらけになって庭作業をしている。
二人がプランターに撒いているのは、黒くて小さなイチゴの種。
\\「このイチゴはね、実がなるまで2年もかかるの」
\\「2ねん？ ながいなぁ」
\\「でもね、時間をかけた分だけ、とびきり甘くなるのよ。真っ赤になったら、二人でお腹いっぱい食べようね」\\
母の言葉に、ノアは目を輝かせて小指を差し出す。
\\「うん！ やくそく！」\\
指切りをする二人の笑顔。幸せの絶頂。

\subsection{シーン２：灰色の寝室}
現在。%画面の彩度は低く、寒色寄りのグレー。
エレナは病に伏せ、痩せ細っている。激しい咳の音だけが響く。
10歳になったノアは、濡れタオルを変え、背中をさするが、その手は微かに震えている。
窓の外、プランターのイチゴにはまだ実がない。緑色の葉が風に揺れているだけだ。

\subsection{シーン３：賢者の庵}
鬱蒼とした森の奥。ノアは息を切らして駆け込んでくる。
\\「賢者様！ 前の薬草、もう効かないんだ。もっと凄いのをちょうだい！ このままじゃお母さんが死んじゃう！」\\
賢者は静かに首を横に振る。
\\「あれ以上の薬はない。あとは母の生命力に任せるしか……」
\\「嫌！ 絶対に嫌！ 何かあるはずだ、賢者様なら知ってるはずだろ！？」\\
ノアは賢者の服を掴み、涙ながらに懇願する。その必死さは、ただの駄々っ子のものではない。
ノアの必死さに、賢者は未来（母が死に、ノアが後を追う絶望的な未来）を予見してしまう。
賢者はため息と共に奥から漆黒のスクロールを取り出す。
\\「……一つだけ、方法がある。だが、これは治療ではない。『取引』だ」
\\「とりひき？」\\
賢者はスクロールを見せる。%代償（存在の消滅と忘却）について語る。
\\「お前が消えれば母は助かる。だが、お前は誰の記憶にも残らない」\\
ノアは言葉を失い、震える。
賢者はスクロールを棚に戻す。
\\「今の話は忘れなさい。……死ぬより辛い選択など、するものではない」

\subsection{シーン４：孤独の天秤}
その夜。ノアはベッドの上で眠れずに膝を抱えている。
%難しいことはわからない。でも、自分が消えることへの生理的な恐怖で身震いする。
\\「消えるって、どうなるんだろう。痛いのかな。暗いのかな」\\
死への根源的な恐怖。生理的な死への恐怖に身震いする。しかし、壁の向こうから母の苦しむ咳が聞こえ、思考が切り替わる。
もし母が死んだら、この冷たい家に自分は一人ぼっちになる。
\\「……ひとりの方が、怖いよ」\\
ノアの中で、「死への恐怖」よりも「母のいない孤独への恐怖」の天秤が重くなる。

\subsection{シーン５：残酷な願い}
翌朝。窓辺のプランターには、イチゴの可愛らしい「白い花」が咲いている。
エレナが目を覚まし、その花を見て微笑む。
\\「……あぁ、見て、ノア。花が咲いたわ」
\\「うん。もうすぐ実になるよ」
\\「楽しみね……。真っ赤になったら、約束通り、二人でお腹いっぱい食べましょうね。……テーブルに並べきれないくらい、たくさん」\\
母の「二人で」という言葉が、鋭い棘のようにノアの胸に刺さる。
賢者の言葉が脳裏をよぎる。『お前が消えれば母は助かる。だが、お前は誰の記憶にも残らない』――それを選べば、この約束は永遠に果たされない。
\\「……うん、そうだね」\\
ノアは声を絞り出すのが精一杯だ。恐怖で足がすくみ、まだ決断できない。
エレナはそんなノアの手を弱々しく、けれど愛おしそうに握る。
\\「お願い、もっと自分を大事にして。あなたが笑ってくれることが、お母さんの一番の幸せなんだから」\\
その言葉は、ノアにとってあまりに残酷なパラドックスとなる。
母を助けたい。でも、助けるためには「自分を大事にして」という母の願いを裏切り、自分を消さなければならない。
「生きたい」という本能と、「母を救いたい」という愛の天秤が、ギリギリで揺れ動く。

\subsection{シーン６：決意}
嵐の夜。エレナが危篤状態となる。ノアは豪雨の中を走り、賢者の小屋へ飛び込む。
\\「お願い！ 貸して！ 今すぐじゃないと間に合わない！」\\
賢者は悲痛な顔で問う。
\\「本気か？ お前自身が消えてしまうのだぞ」\\
ノアは泣きじゃくりながら、子供なりの限界を叫ぶ。
\\「怖いよ！ ……でも、お母さんがいなくなる方がもっと怖いんだ！」\\
その叫びは、言葉よりも遥かに切実な10歳の少年の本音だった。
賢者は目を閉じ、その慟哭を受け止める。これほど純粋な魂の叫びを、大人の理屈で止めることはできない。
\\「……行け。その愛、誰にも止められまい」\\
賢者は震える手でスクロールを差し出す。

\subsection{シーン７：消失}
戻ってきたノアは、母のベッドサイドで魔法陣を開く。
自身の体が光の粒子となって崩れ始める中、ノアは窓の外のイチゴへ視線を投げる。
\\「ごめんね、お母さん。約束、守れそうにないや」\\
光が強くなる。ノアの表情に恐怖はない。
\\「……よかった。これで、ひとりじゃないね」
\\「ありがとう……お母さん……」\\
光が弾け、画面がホワイトアウトする。

\subsection{シーン８：歴史の改変}
\textbf{（重要演出）}世界が再構築される映像。
冒頭の「二人で苗を植える回想シーン」がフラッシュバックする。
しかし、ノイズと共にノアの姿だけがフッと消滅し、「エレナがたった一人で、汗を拭いながら苗を植えている映像」へと上書きされる。
現実の部屋でも、ノアの服やおもちゃが、砂のようにサラサラと崩れて消滅していく。
ベッドのエレナの顔色はみるみる良くなり、穏やかな寝息を立て始める。

\subsection{シーン９：綻びのある日常}
数ヶ月後。晴天。庭には真っ赤なイチゴ。
元気になったエレナが、庭のテーブルに紅茶を置く。
カチャン、カチャン。
無意識に「二つ」のカップを置き、向かいの席に誰もいないことに気づいて、不思議そうに首を傾げる。
\\「……あら？ 私ったら、どうして二つ淹れたのかしら」\\
彼女は胸をさする。病気は治ったはずなのに、胸の奥に、風が吹き抜けるような空洞がある。

\subsection{シーン１０：賢者の来訪}
エレナは山盛りのイチゴが入ったカゴを抱えている。
庭でイチゴを収穫していたエレナが、通りかかった賢者に気づく。
\\「あら、賢者様。森から出てこられるなんて珍しいですね」
\\「ふむ、ちと散歩にな。……豊作ですね」\\
賢者がイチゴに目を向ける。エレナは笑顔でカゴを差し出す。
\\「ええ。でも不思議なんです。私ひとりじゃ、どうしても食べきれなくて……。」\\
賢者はその赤く輝く果実を見つめる。それは誰かの命そのものであり、エレナへの愛の結晶だ。賢者の目に、安堵と哀しみが混ざる。
賢者は首を横に振り、優しく断る。
\\「いいえ。それはあなたが食べるべきだ」
\\「え？」
\\「その赤さは、あなたが注いだ愛情と……長い年月の証です。あなたがその甘さを噛み締めてやれば、胸のつかえも取れるだろう」
\\「はぁ……そういうものでしょうか」

\subsection{シーン１１：涙の味}
エレナは不思議そうにしながらも、賢者の言葉に従い、一粒のイチゴを口に含む。
\\「……おいしい」\\
その瞬間、ボロボロと大粒の涙がこぼれ落ちる。
\\「おかしいわね。すごく甘くて……すごく、優しい味がするの」\\
泣き崩れながら、イチゴを大切そうに抱きしめるエレナ。
脳は忘れていても、身体が失われた半身を覚えている。
賢者はその姿に背を向け、歩き出しながら、誰にも聞こえない声で呟く。
\\「……あぁ、見事だ」\\
虚空へ向けて、慈しむように。
\\「お前の愛は、ちゃんと届いておるよ。……ノア」\\
風が吹き抜け、イチゴの葉が揺れる。%カメラが青空へとパンアップして終幕。

\textbf{「完」}



\clearpage\section{設定}
\begin{description}
    \item [ジャンル] 異世界ファンタジー、感動ドラマ
    \item [タイトル] 約束の果実
    \item [ノア（10歳）]
    \begin{itemize}
        \item 主人公。父を早くに亡くし、母を守ろうとする健気な子供。
        \item \textbf{【性別不詳】} あえて性別は明言せず、中性的な容姿（ショートヘアなど）とする。観客が自身の子供や自分自身を投影できるようにするため。
        \item 「自分の死」よりも「孤独」を恐れる、子供らしい切実な動機を持つ。
        \item 最終的に、母の「自分を大事にして」という願いを裏切ってでも、母の命を救うことを選ぶ（愛ある裏切り）。
    \end{itemize}
    \item [エレナ（30代）]
    \begin{itemize}
        \item ノアの母。不治の病に侵される。
        \item 自分の命よりも息子の幸福を願う、無償の愛の持ち主。
        \item 魔法の代償により、ノアに関する一切の記憶（妊娠、出産、育児）を失う。
    \end{itemize}
    \item [賢者（老齢）]
    \begin{itemize}
        \item 森の奥に住む観測者。世界の理を知る存在。世界の均衡と秩序を守る番人。
        \item ノアの決断を見届け、最後にエレナへ「真実（愛）」を伝える役割を担う。
    \end{itemize}

    \item [重要アイテム]
    \begin{itemize}
        \item \textbf{禁忌のスクロール}
        \begin{itemize}
            \item 賢者が管理している、漆黒の古びた巻物。
            \item \textbf{効果：} いかなる病や呪いも打ち消す奇跡を起こす。
            \item \textbf{代償：} 「等価交換」ではなく、術者の「存在そのもの」を燃料とする。使用者は光となって消滅し、世界中の人々の記憶から抹消されるだけでなく、写真や私物などの「生きた痕跡」も歴史から排除（修正）される。
        \end{itemize}
        \item \textbf{イチゴ（約束の果実）}
        \begin{itemize}
            \item \textbf{種（過去）：} 2年前、母と二人で撒いた小さな希望。
            \item \textbf{緑/白（未熟）：} 成長途中のノア、未完成の約束。
            \item \textbf{赤（完熟）：} ノアの命、母への愛、完成した約束。
            \item この物語ではイチゴを単なる果物ではなく、「ノアの化身」として扱います。
        \end{itemize}
    \end{itemize}

\end{description}


\clearpage\section{キャラクター心理}

\subsection{ノアの感情曲線}
本作品の主人公ノアは、短い尺の中で劇的な心理変化を遂げます。%演出および作画においては、以下の4段階の感情推移を意識してください。

\subsubsection{第１段階：恐怖と依存（シーン２〜３）}
\textbf{「お母さんが死ぬのが怖い。自分が消えるのも怖い」}
\begin{itemize}
    \item 賢者の前では必死ですが、一人になると年相応の弱さを見せます。
    \item 恐怖のあまり、無意識に自分の腕を抱いたり、震えたりする描写を入れます。まだ「覚悟」はできていません。
\end{itemize}

\subsubsection{第２段階：孤独の天秤（シーン４）}
\textbf{「死ぬこと（肉体の消滅）」vs「孤独（精神の死）」}
\begin{itemize}
    \item ここが最大の転換点です。大人びた自己犠牲ではなく、\textbf{「一人ぼっちは嫌だ」という子供の本能的な恐怖}が、死への恐怖を上回る瞬間です。
    \item ヒーローとしての決断ではなく、「逃避（孤独からの逃げ）」としての側面を持たせることで、より切実さを演出します。
\end{itemize}

\subsubsection{第３段階：皮肉な決意（シーン５）}
\textbf{「お母さんの願いを裏切ることが、お母さんを守る唯一の道」}
\begin{itemize}
    \item 母から「自分を大事にして」と言われた時、ノアは泣きそうな顔で笑います。
    \item この瞬間、彼は「子供」から「守る者（親のような視点）」へと精神的に成長します。
\end{itemize}

\subsubsection{第４段階：安寧と昇華（シーン６〜７）}
\textbf{「これで、もう寂しくない」}
\begin{itemize}
    \item 魔法を使う瞬間、ノアに悲壮感はありません。むしろ、母を救える安堵と、孤独から解放される喜びで、穏やかな表情（聖母のような微笑み）をしています。
    \item \textbf{「消えること＝救い」}として描くことで、観客の涙を誘います。
\end{itemize}

\subsection{賢者の行動原理について} 
なぜ賢者は子供に禁忌の品を渡したのか。
それは彼がノアの未来を予見し、「母の死後、ノアが絶望して後を追う（共倒れになる）未来」を感じ取ったからです。
賢者はノアを子供扱いせず、その凄絶な愛の深さを認め、「どう生きるか（どう終わるか）」の選択権を彼に委ねました。
これは残酷な判断ですが、賢者なりの最大限の慈悲（救済措置）として描きます。

\clearpage\section{各シーンの演出・脚本意図解説}
\subsection{シーン１〜２：色彩による対比}
\textbf{［狙い］「幸せな過去」と「絶望的な現在」の可視化} \\
冒頭（過去）は暖色系のフィルターで「楽園」のように描き、現在シーンでは彩度を落とした寒色（死の色）で描きます。このギャップにより、ノアが取り戻したいものが何であるかを一目で理解させます。

\subsection{シーン４：孤独の天秤（動機の強化）}
\textbf{［狙い］子供視点での「自己犠牲」のリアリティ} \\
子供がいきなり「母のために死ねる」と決意するのは嘘くさくなりがちです。
そこで、「死ぬのは怖い」という本音を吐露させた上で、「でも独りぼっちになるのはもっと怖い」という天秤を描きます。これにより、彼の行動原理に強い説得力を持たせます。

\subsection{シーン５：残酷な願い（皮肉）}
\textbf{［狙い］「愛」が「自己犠牲」のトリガーになる矛盾} \\
母の「自分を大事にして」という言葉が、逆にノアの背中を押してしまう（母を悲しませないために、自分が消えるしかない）という皮肉な構造を作ります。この「すれ違う愛」が物語の悲劇性を高めます。

\subsection{シーン８：歴史の改変（映像的クライマックス）}
\textbf{［狙い］「消える」ことの残酷さの映像化} \\
単にノアが光になって消えるだけでは不十分です。「過去の回想シーン」にノイズが走り、ノアの姿だけが消去され、最初から母一人だったかのように映像が書き換わる演出を入れます。
これにより、ノアが「死んだ」のではなく「存在しなかったことになった」という絶望的なルールを、視覚的衝撃と共に伝えます。

\subsection{シーン１１：味覚による記憶の呼び覚まし}
\textbf{［狙い］理屈を超えた「魂の記憶」} \\
脳（記憶）はノアを忘れていますが、身体（味覚・涙）は彼を覚えている。
「おいしい」というシンプルな感覚から涙が溢れる描写にすることで、言葉で説明するよりも雄弁に「愛が残ったこと」を証明します。
最後に賢者が名前を呼ぶことで、観客が抱える「誰か覚えていてあげて！」という感情を救済し、温かい余韻で幕を閉じます。

\clearpage\section{物語の伏線とテーマの深化}

\subsection{「2年」の時間経過（構造の伏線）}
\begin{itemize}
    \item \textbf{伏線：} 冒頭で「種からだと実がなるまで2年かかる」と提示し、中盤で「もうすぐ2年だ」とノアがつぶやきます。蕾はまだ白く、時間（命）が足りないことを示唆します。
    \item \textbf{回収：} ラストシーンでイチゴは真っ赤に熟しています。これは、ノアが本来生きるはずだった未来の時間や、彼の命そのものが、魔法によって一瞬でイチゴに注ぎ込まれ、空白の時間を埋めたことを意味します。
\end{itemize}

\subsection{「一人で種を撒いた」という矛盾（映像の伏線）}
\begin{itemize}
    \item \textbf{伏線：} 冒頭では「二人」で種を撒いていました。
    \item \textbf{回収：} 歴史改変後、回想シーンが「エレナが一人で種を撒いている映像」に書き換わります。エレナは「私一人じゃ食べきれない」「一人で持て余している」と繰り返します。これは事実として「一人で育てた歴史」になったからですが、観客には「そこにいたはずのノア」の不在を痛烈に感じさせるフックとなります。
\end{itemize}

\subsection{賢者の立ち位置（視点の伏線）}
\begin{itemize}
\item \textbf{伏線：} 賢者は決して魔法を推奨せず、あくまで「観測者」として振る舞います。
\item \textbf{回収：} ラストで彼だけがノアの名前を呼ぶことで、彼が単なる魔法使いではなく、「この悲しくも美しい奇跡の唯一の証人」であったことが確定します。
\end{itemize}

\clearpage\section{感動を生むメカニズム（涙へのプロセス）}
本企画は「悲しいから泣ける」ではなく、「愛の深さに触れて泣ける」構造を目指します。

\begin{enumerate}
\item \textbf{不可逆性の提示と予感}（シーン3〜4）
\begin{itemize}
\item 仕掛け： 「存在の消滅」という重すぎる代償と、それでも母を想う子供の健気さを描きます。
\item 観客の心理： 「助かってほしい」という願いと、「でも消えてしまう」という予感の間でサスペンス（緊張感）を高めます。
\end{itemize}
\item \textbf{歴史改変による喪失の実感}（シーン7〜9）
\begin{itemize}
\item 仕掛け： 幸せな回想シーンが「孤独な映像」に上書きされる残酷な演出。
\item 観客の心理： ノアが「いなくなった」ことの決定的な証拠を突きつけられ、猛烈な喪失感（悲しみ）を与えます。
\end{itemize}
\item \textbf{感覚的な再会と救済}（シーン10〜11）
\begin{itemize}
\item 仕掛け： 記憶はないのに涙が止まらない母と、それを肯定する賢者の言葉。
\item 観客の心理：
\begin{itemize}
\item 「頭では忘れても、心は覚えている」という現象により、ノアの犠牲が報われたことを感じさせます。
\item 賢者が「見事だ」と称賛し名前を呼ぶことで、観客の「ノアの存在を認めてほしい」という欲求を満たし、悲しみを「温かい感動」へと昇華させます。
\end{itemize}
\end{itemize}
\end{enumerate}



\clearpage\section{企画意図・解説資料}
\subsection{物語の核：「矛盾」と「昇華」の設計}
この物語は、単なるお涙頂戴ではなく、人間の複雑な心理を描くために以下の構成にしています。
\begin{description}
    \item [動機のリアリティ（矛盾）] 
    主人公ノアの決断を、美しい自己犠牲としてだけでなく、子供特有の\textbf{「死ぬのは怖いが、独りぼっちになるのはもっと怖い」}という切実なエゴ（生存本能の裏返し）として描きます。これにより、観客は「聖人君子」としてではなく「等身大の子供」としてノアに感情移入します。
    \item [愛の裏切り（昇華）] 
    「自分を大事にして」という母の願いを叶えるためには、母を見殺しにするしかない。逆に、母を救うためには、その願いを裏切って自分が消えるしかない。この\textbf{「愛ゆえに親の願いを裏切る」}という構造を作ることで、ラストシーンの切なさを最大化します。
\end{description}

\iffalse
\subsection{視覚的演出：「色」と「記憶」の対比}
セリフで説明するのではなく、画面の色彩設計そのものをストーリーテリングの一部として機能させます。
\begin{itemize}
    \item 「\textbf{彩度の低い世界（灰色）}」 ＝ 現在、死の予感、閉塞感
    \item 「\textbf{鮮やかな赤（イチゴ）}」 ＝ ノアの命、情熱、残された愛
\end{itemize}
物語の前半を意図的に彩度を落として描き、ラストシーンの青空と「イチゴの赤」を強烈なコントラストで見せることで、ノアの命が世界に色を取り戻したことを視覚的に表現します。
また、歴史改変シーンにおける「ノイズによる人物消去」の演出は、魔法の残酷さをデジタル的な違和感として突きつけ、観客の心に深い爪痕を残す狙いがあります。
\fi

\subsection{キャラクター戦略：性別の意図的な排除}
主人公ノアの性別をあえて設定せず、中性的なデザインと演技（一人称「ボク」、高めの声）で描写します。
これにより、以下の効果を狙います。
\begin{itemize}
    \item \textbf{感情移入の最大化：} 観客はノアを「男の子」や「女の子」としてではなく、普遍的な\textbf{「我が子」という概念}として認識するため、観客の没入感を高めます。
    \item \textbf{テーマの純化：} ジェンダーによる役割（男だから強くあれ、女だから優しくあれ等）をノイズとして排除し、純粋な「親を想う子供の魂」だけを抽出して描きます。
\end{itemize}

\subsection{観客へのアプローチ}
過去の制作物に対するアンケート結果を分析し、弱点を克服するための構成にしました。

\begin{description}
    \item [課題１：「結末の消化不良感」] 
    \begin{description}
        \item []
        \item [アンケートの声] 『中途半端なところで終わってしまった感があった』『見終わったあとの「なんだあれは」というモヤモヤ感』『結末がどうなったのか分かりづらかった』『2話目ください（＝1話で完結していない不満）』
        \item [本作の解決策] 
        「ノアが消えて終わり（バッドエンド）」ではなく、その数ヶ月後の「母が救われた姿（救済）」までをしっかりと描きます。さらに賢者が「愛は届いている」と明言することで、物語としての「答え」を提示し、心地よい余韻と共に完結させます。
    \end{description}
    \item [課題２：「音声・セリフの聞き取りづらさ」] 
    \begin{description}
        \item []
        \item [アンケートの声] 『声優の声が小さく内容が入ってこなかった』『ちょっと聞き取れない部分がありました』
        \item [本作の解決策] 
            \begin{itemize}
                \item \textbf{「静寂」の活用:} 本作は叫び合うシーンがほとんどありません。BGMをカットし、環境音（雨音、衣擦れ）だけにするシーンを作ることで、ボソボソとした呟きでも際立つ音響設計にします。
                \item \textbf{視覚情報の補完:} 賢者の最後のセリフなど、聞こえにくくても文脈で伝わる演出（口元のアップや表情）を多用し、聴覚だけに頼らない設計にします。
            \end{itemize}
    \end{description}
    \item [課題３：「ドラマパートの物足りなさ」] 
    \begin{description}
        \item []
        \item [アンケートの声] 『日常パートの書き込み量が戦闘シーンと比べて少ない』『教室での会話で、真正面の会話が続いた』
        \item [本作の解決策] 
        本作は派手なアクションがない分、作画のリソースを「芝居（演技）」に全振りします。「震える指先」「視線の泳ぎ」「イチゴを食べる口元」といった\textbf{微細な演技}を丁寧に書き込むことで、静かな会話劇であっても画面を持たせ、退屈させないドラマを作ります。
    \end{description}
\end{description}


\iffalse
\clearpage\section{制作上の懸念点と対策}
\subsection{歴史改変（シーン7）の分かりにくさ懸念}
「過去が書き換わる」というSF的な現象を、セリフでの説明なしに理解させる必要があります。

\textbf{【対策：明確なグリッチ演出】}
回想シーンにデジタルノイズのようなグリッチ（歪み）を入れ、ノアの姿が「切り抜かれる」ように消え、背景が自然に繋がるようなアニメーション処理を施します。これにより「最初からいなかったことになった」という違和感を強調します。

\subsection{賢者のラストのセリフの演技}
賢者の最後のセリフが説明的になりすぎると、余韻を壊す恐れがあります。

\textbf{【対策：ささやき声と環境音】}
賢者の「……ノア」という最後の言葉は、エレナには聞こえないレベルの吐息混じりの演技にします。直後に風の音（SE）とBGMの盛り上がりを重ねることで、言葉の意味よりも「想いが風に乗って空へ届く」というニュアンスを重視した音響演出を行います。
\fi
